\documentclass[11pt,a4paper]{article}

\usepackage[T1]{fontenc}
\usepackage[utf8]{inputenc}
\usepackage[ngerman]{babel}
\usepackage{lmodern}
\usepackage{microtype}
\usepackage[a4paper,margin=2.5cm]{geometry}
\usepackage{amsmath,amssymb,amsthm}
\usepackage{mathtools}
\usepackage{hyperref}

\hypersetup{
	colorlinks=true,
	linkcolor=blue,
	citecolor=blue,
	urlcolor=blue,
	pdftitle={Die Zwölf-Linie bei ungeraden negativen Zetawerten},
	pdfauthor={Thomas Hoffbauer und ChatGPT},
	pdfsubject={Bernoulli-Zahlen, Zetafunktion, Nennerstruktur},
	pdfkeywords={Riemannsche Zetafunktion, Bernoulli-Zahlen, von Staudt-Clausen, negative Spezialwerte}
}

\title{\textbf{Die Zwölf-Linie bei ungeraden negativen Zetawerten}\\
	\large Eine arithmetische Beobachtung zur Nennerstruktur von $\zeta(-(2m-1))$}
\author{Thomas Hoffbauer \and ChatGPT}
\date{\today}

\theoremstyle{definition}
\newtheorem{definition}{Definition}[section]
\newtheorem{beispiel}[definition]{Beispiel}

\theoremstyle{plain}
\newtheorem{satz}[definition]{Satz}
\newtheorem{proposition}[definition]{Proposition}
\newtheorem{lemma}[definition]{Lemma}
\newtheorem{korollar}[definition]{Korollar}
\newtheorem{vermutung}[definition]{Vermutung}
\newtheorem{empirischebehauptung}[definition]{Empirische Behauptung}

\theoremstyle{remark}
\newtheorem{bemerkung}[definition]{Bemerkung}

\begin{document}
	\maketitle
	
	\begin{abstract}
		Wir untersuchen die vollständig gekürzten Nenner der ungeraden negativen Spezialwerte der Riemannschen Zetafunktion
		\[
		\zeta(-(2m-1))=-\frac{B_{2m}}{2m}.
		\]
		Dabei zeigt sich numerisch eine auffällige Teilmenge von Indizes $m$, für die der Endnenner genau gleich $12$ ist. Diese Teilmenge wird hier als \emph{Zwölf-Linie} bezeichnet. Ausgehend von der von-Staudt--Clausenschen Nennerstruktur der Bernoulli-Zahlen formulieren wir eine exakte Teilbarkeitscharakterisierung der Zwölf-Indizes und leiten daraus eine Reihe empirischer Vermutungen über Kongruenzklassen und sogenannte Fast-Zwölf-Fälle ab.
	\end{abstract}
	
	\section{Einleitung}
	
	Die negativen ganzzahligen Spezialwerte der Riemannschen Zetafunktion sind durch die Bernoulli-Zahlen gegeben:
	\[
	\zeta(-n)=-\frac{B_{n+1}}{n+1}.
	\]
	Für ungerade negative Stellen erhält man insbesondere
	\[
	\zeta(-(2m-1))=-\frac{B_{2m}}{2m}
	\qquad (m\ge 1).
	\]
	Da die Bernoulli-Zahlen rational sind, sind auch diese Zetawerte rationale Zahlen. Die vorliegende Notiz befasst sich nicht mit ihren Größenordnungen, sondern mit ihrer \emph{Nennerstruktur nach vollständiger Kürzung}.
	
	Numerisch fällt auf, dass der Nenner $12$ ungewöhnlich häufig als Endnenner auftritt. Der prototypische Fall ist
	\[
	\zeta(-1)=-\frac{1}{12},
	\]
	doch auch bei wesentlich größeren Indizes tritt derselbe Endnenner erneut auf. Dies legt nahe, die entsprechende Teilmenge systematisch zu untersuchen.
	
	\section{Grundstruktur}
	
	Schreiben wir die Bernoulli-Zahl $B_{2m}$ in vollständig gekürzter Form als
	\[
	B_{2m}=\frac{A_{2m}}{D_{2m}},
	\]
	wobei
	\[
	A_{2m}\in\mathbb{Z},\qquad D_{2m}\in\mathbb{N},\qquad \gcd(A_{2m},D_{2m})=1,
	\]
	so folgt
	\[
	\zeta(-(2m-1))=-\frac{A_{2m}}{(2m)D_{2m}}.
	\]
	
	Der Satz von von Staudt--Clausen liefert den Nenner von $B_{2m}$ explizit:
	\[
	D_{2m}=\prod_{p-1\mid 2m} p,
	\]
	wobei das Produkt über alle Primzahlen $p$ läuft, deren Vorgänger $p-1$ ein Teiler von $2m$ ist.
	
	Damit ist der \emph{Rohnenner} des Zetawertes
	\[
	R_m:=(2m)D_{2m}.
	\]
	
	\begin{definition}
		Sei
		\[
		\zeta(-(2m-1))=\frac{P_m}{Q_m}
		\]
		in vollständig gekürzter Form mit
		\[
		P_m\in\mathbb{Z},\qquad Q_m\in\mathbb{N},\qquad \gcd(P_m,Q_m)=1.
		\]
		Dann nennen wir:
		\begin{itemize}
			\item $m$ einen \emph{Zwölf-Index}, falls $Q_m=12$,
			\item $m$ einen \emph{Fast-Zwölf-Index}, falls $Q_m=12r$ mit einem ganzzahligen $r>1$.
		\end{itemize}
	\end{definition}
	
	\section{Exakte Charakterisierung}
	
	\begin{proposition}
		Mit den obigen Bezeichnungen gilt:
		\[
		Q_m=12
		\quad\Longleftrightarrow\quad
		\frac{(2m)D_{2m}}{12}\mid A_{2m}.
		\]
	\end{proposition}
	
	\begin{proof}
		Da
		\[
		\zeta(-(2m-1))=-\frac{A_{2m}}{(2m)D_{2m}},
		\]
		ist der vollständig gekürzte Nenner
		\[
		Q_m=\frac{(2m)D_{2m}}{\gcd(A_{2m},(2m)D_{2m})}.
		\]
		Also gilt $Q_m=12$ genau dann, wenn
		\[
		\gcd(A_{2m},(2m)D_{2m})=\frac{(2m)D_{2m}}{12}.
		\]
		Dies ist äquivalent dazu, dass
		\[
		\frac{(2m)D_{2m}}{12}\mid A_{2m}.
		\]
	\end{proof}
	
	\begin{bemerkung}
		Die Aussage ist exakt und nicht heuristisch. Sie zeigt, dass die Zwölf-Linie vollständig durch eine Teilbarkeitsbedingung an den Bernoulli-Zähler $A_{2m}$ beschrieben wird.
	\end{bemerkung}
	
	\section{Die Zwölf-Linie als arithmetische Struktur}
	
	Im numerisch untersuchten Bereich zeigt sich, dass Zwölf-Indizes nur in bestimmten Kongruenzklassen auftreten.
	
	\begin{empirischebehauptung}
		Für alle numerisch untersuchten Zwölf-Indizes $m\le 100$ gilt
		\[
		m\equiv 1 \text{ oder } 5 \pmod 6.
		\]
		Äquivalent dazu erfüllen die zugehörigen Bernoulli-Indizes
		\[
		2m\equiv 2 \text{ oder } 10 \pmod{12}.
		\]
	\end{empirischebehauptung}
	
	\begin{bemerkung}
		Diese Beobachtung ist zunächst rein empirisch. Sie bedeutet, dass in dem untersuchten Bereich der Endnenner $12$ nur dann auftritt, wenn $m$ ungerade ist und nicht durch $3$ teilbar wird.
	\end{bemerkung}
	
	Noch stärker scheint die Klasse
	\[
	m\equiv 1 \pmod 6
	\]
	bevorzugt zu sein.
	
	\begin{vermutung}[Zwölf-Linien-Vermutung]
		Die Menge
		\[
		\mathcal{Z}_{12}:=\{m\in\mathbb{N}: Q_m=12\}
		\]
		bildet eine arithmetisch dünne, aber strukturierte Teilmenge der natürlichen Zahlen. Numerisch deutet alles darauf hin, dass Zwölf-Indizes stark durch die Kongruenzbedingung
		\[
		m\equiv 1 \text{ oder } 5 \pmod 6
		\]
		vorgefiltert werden und innerhalb dieser beiden Klassen die Restklasse
		\[
		m\equiv 1 \pmod 6
		\]
		deutlich bevorzugen.
	\end{vermutung}
	
	\section{Fast-Zwölf-Fälle}
	
	Für Fast-Zwölf-Indizes ist der Endnenner von der Form
	\[
	Q_m=12r
	\]
	mit einem kleinen Restfaktor $r>1$. Diese Fälle sind besonders aufschlussreich, da sie oft nur um einen einzigen nicht absorbierten Primfaktor vom echten Zwölf-Fall entfernt sind.
	
	\begin{definition}
		Sei
		\[
		R_m=(2m)D_{2m}=12K_m.
		\]
		Dann nennen wir
		\[
		r_m:=\frac{Q_m}{12}
		\]
		den \emph{Restfaktor} des Index $m$.
	\end{definition}
	
	\begin{bemerkung}
		Ist $m$ ein Zwölf-Index, so gilt $r_m=1$. Ist $m$ ein Fast-Zwölf-Index, so misst $r_m$ genau den Anteil von $K_m$, der durch den Bernoulli-Zähler $A_{2m}$ \emph{nicht} absorbiert wird.
	\end{bemerkung}
	
	\begin{vermutung}[Fast-Zwölf-Phänomen]
		Fast-Zwölf-Indizes entstehen typischerweise dadurch, dass im Rohnenner
		\[
		(2m)D_{2m}=12K_m
		\]
		nicht alle Primfaktoren von $K_m$ im Bernoulli-Zähler $A_{2m}$ erscheinen. In vielen numerischen Beispielen bleibt nur ein einzelner kleiner Primfaktor oder ein kleines Primprodukt übrig.
	\end{vermutung}
	
	\section{Von-Staudt-induzierte Restfaktoren}
	
	Die Restfaktoren der Fast-Zwölf-Fälle scheinen oft direkt mit der von-Staudt--Clausenschen Struktur zusammenzuhängen. Wenn nämlich eine Primzahl $p$ die Bedingung
	\[
	p-1\mid 2m
	\]
	erfüllt, tritt sie automatisch im Nenner $D_{2m}$ auf. Entscheidend ist dann, ob derselbe Faktor auch den Bernoulli-Zähler $A_{2m}$ teilt.
	
	\begin{bemerkung}
		Numerisch treten insbesondere Faktoren wie $11$ und $23$ gehäuft als Restfaktoren auf. Dies ist heuristisch dadurch erklärbar, dass sie über die Bedingungen
		\[
		10\mid 2m \qquad\text{bzw.}\qquad 22\mid 2m
		\]
		natürlich in den Nenner $D_{2m}$ eingespeist werden. Die zugehörigen Fast-Zwölf-Fälle können daher als solche Indizes interpretiert werden, bei denen genau ein strukturell erzeugter Nennerfaktor der vollständigen Kürzung widersteht.
	\end{bemerkung}
	
	\section{Der prototypische Fall \texorpdfstring{$\zeta(-1)=-1/12$}{zeta(-1)=-1/12}}
	
	Der Wert
	\[
	\zeta(-1)=-\frac{1}{12}
	\]
	nimmt innerhalb der Zwölf-Linie eine Sonderstellung ein. Einerseits ist er der historisch bekannteste regulierte Spezialwert der Zetafunktion. Andererseits ist er zugleich der elementarste Zwölf-Fall überhaupt.
	
	Denn aus
	\[
	B_2=\frac16
	\]
	folgt sofort
	\[
	\zeta(-1)=-\frac{B_2}{2}=-\frac{1}{12}.
	\]
	Hier entsteht der Nenner $12$ unmittelbar aus der Bernoulli-Struktur, ohne dass eine komplizierte zusätzliche Kürzung erforderlich wäre. In diesem Sinn kann $\zeta(-1)=-1/12$ als \emph{Urform} der gesamten Zwölf-Linie angesehen werden.
	
	\section{Ausblick}
	
	Die hier formulierten Beobachtungen legen mehrere weiterführende Fragestellungen nahe:
	
	\begin{itemize}
		\item Besitzt die Zwölf-Linie eine asymptotisch beschreibbare Dichte innerhalb geeigneter Kongruenzklassen?
		\item Lässt sich die Dominanz der Klasse $m\equiv 1\pmod 6$ theoretisch erklären?
		\item Welche Rolle spielen reguläre und irreguläre Primzahlen für das Auftreten von Fast-Zwölf-Fällen?
		\item Kann die Restfaktor-Landschaft systematisch über die Teilbarkeit der Bernoulli-Zähler beschrieben werden?
	\end{itemize}
	
	Die Zwölf-Linie erscheint damit als ein kleines, aber bemerkenswert strukturiertes Fenster in die Arithmetik der Bernoulli-Zahlen und der negativen Spezialwerte der Zetafunktion.
	
\end{document}